\documentclass[10pt, letterpaper]{article}

\usepackage[margin=0.72in]{geometry}
\usepackage{array}
\usepackage{booktabs}
\usepackage{longtable}
\usepackage{tabularx}
\usepackage{pdflscape}
\usepackage[table]{xcolor}
\usepackage{titlesec}
\usepackage{fancyhdr}
\usepackage{enumitem}
\usepackage{microtype}
\usepackage[hidelinks]{hyperref}

\definecolor{headingblue}{HTML}{1F4E79}
\definecolor{lightblue}{HTML}{D9EAF7}
\definecolor{softgray}{HTML}{F2F4F7}
\definecolor{linegray}{HTML}{A6A6A6}

\pagestyle{fancy}
\fancyhf{}
\lhead{\small Md Kamruzzaman Kamrul}
\rhead{\small EDCI 59100-016 \textbar{} Summer 2026}
\cfoot{\small \thepage}
\renewcommand{\headrulewidth}{0.4pt}
\setlength{\headheight}{14pt}

\titleformat{\section}{\normalfont\large\bfseries\color{headingblue}}{\thesection.}{0.45em}{}
\titleformat{\subsection}{\normalfont\normalsize\bfseries\color{headingblue}}{\thesubsection.}{0.45em}{}
\titlespacing*{\section}{0pt}{10pt}{4pt}
\titlespacing*{\subsection}{0pt}{7pt}{3pt}

\setlength{\parindent}{0pt}
\setlength{\parskip}{4pt}
\setlist[itemize]{leftmargin=*, itemsep=2pt, topsep=2pt}
\setlist[enumerate]{leftmargin=*, itemsep=2pt, topsep=2pt}
\renewcommand{\arraystretch}{1.16}

\newcolumntype{L}[1]{>{\raggedright\arraybackslash}p{#1}}
\newcolumntype{Y}{>{\raggedright\arraybackslash}X}

\begin{document}

\begin{center}
{\Large\bfseries Synthesis Matrix and Organizational Outline}\\[0.35em]
{\normalsize Draft Literature Review Plan}\\[0.25em]
{\normalsize Md Kamruzzaman Kamrul}\\
{\normalsize EDCI 59100-016 \textbar{} Summer 2026}
\end{center}

\vspace{0.6em}

\section{Purpose and Submission Map}

This assignment moves the review from individual article notes to a working structure for synthesis. The review topic is the use of vision-language models and related multimodal artificial intelligence systems for construction-site safety monitoring and activity tracking.

Spatiotemporal reasoning is used here in a simple way. It means understanding what is happening, where it is happening, and how that situation changes over time. On a construction site, this may include PPE use, unsafe proximity, equipment activity, work progress, or whether a visual scene matches a safety rule.

\textbf{Review type.} The final review is planned as a scoping review using PRISMA-ScR guidance. This fits the project because the literature is still developing across construction management, computer vision, natural language processing, robotics, and safety science. The studies use different data sources, model types, and evaluation metrics. A scoping review can chart that range without forcing the papers into a narrow effect-size comparison.

\section{Source Selection}

The full working corpus includes baseline, transition, core, and review papers. For this assignment, I selected 15 key sources. The goal is not to include every paper. The goal is to choose papers that help reveal patterns, contrasts, and gaps.

The selected set includes:

\begin{itemize}
\item baseline papers that show the earlier dependence on sensors, manual reporting, and computer vision;
\item transition papers that move from detection toward semantic reasoning, captioning, and activity recognition;
\item core VLM or multimodal studies for safety monitoring and progress tracking;
\item review papers that help locate broader gaps in construction AI and computer vision.
\end{itemize}

\section{Synthesis Matrix}

The matrix below builds from the earlier source evaluation notes. It adds theme, comparison, tension, and gap columns. This turns the working notes into a structure for the literature review.

\begin{landscape}
\small
\rowcolors{2}{softgray}{white}
\begin{longtable}{L{3.0cm} L{2.2cm} L{2.5cm} L{2.9cm} L{3.5cm} L{3.6cm} L{3.2cm}}
\caption{Synthesis Matrix for Key Sources}\\
\toprule
\rowcolor{lightblue}
\textbf{Source} & \textbf{Task and evidence type} & \textbf{Method or model} & \textbf{Main finding} & \textbf{Connection to other studies} & \textbf{Tension or contrast} & \textbf{Gap or use in outline} \\
\midrule
\endfirsthead
\toprule
\rowcolor{lightblue}
\textbf{Source} & \textbf{Task and evidence type} & \textbf{Method or model} & \textbf{Main finding} & \textbf{Connection to other studies} & \textbf{Tension or contrast} & \textbf{Gap or use in outline} \\
\midrule
\endhead
\bottomrule
\endfoot

Zhong et al. (2019) & Construction computer vision review. Baseline map. & Bibliometric and scientometric review of 216 Web of Science papers. & Computer vision was already used for inspection, safety monitoring, performance analysis, and tracking before current VLM work. & Establishes the older CV landscape that later VLM papers build on. Connects to Nath et al. and Sherafat et al. as pre-VLM foundations. & The review shows growth in image and video analysis, but most work remains visual and task-specific. & Use in the introduction and Theme 1 to show why the field needed stronger semantic interpretation. \\

Sherafat et al. (2020) & Activity recognition review. Baseline for worker and equipment monitoring. & State-of-the-art review of automated activity recognition methods. & Many systems reached high accuracy for selected worker or equipment tasks, often using sensors, video, or classifiers. & Provides a temporal and activity-recognition base for Jeoung et al. and Xiao et al. & Accuracy is often strong in controlled task settings, but the methods do not fully explain site meaning or link activity to reports. & Use in Theme 3 to show the earlier activity-recognition branch. \\

Nath et al. (2020) & PPE detection and field safety. Baseline empirical study. & Real-time deep learning detection of workers, hardhats, and safety vests. & The study shows that CV can automate PPE checks with strong detection performance. & Connects directly to Gil and Lee because both study PPE compliance. It also shows what later captioning and VLM approaches try to improve. & PPE detection is measurable and useful, but it does not capture broader hazard context. & Use in Theme 2 as the clear baseline for safety compliance monitoring. \\

Wu et al. (2021) & Safety reasoning. Transition paper. & Computer vision combined with semantic reasoning and ontology-based representation. & The paper moves beyond object detection by using logical relationships for hazard identification. & Bridges Nath et al.'s detection approach and later VLM or ontology-supported systems such as Chan et al. & The method adds reasoning, but the reasoning remains external to the visual model. & Use in Theme 1 and Theme 2 to show the move from seeing objects to interpreting safety meaning. \\

Liu et al. (2020) & Construction activity scene description. Transition paper. & Image captioning models evaluated with language generation metrics. & The study shows that construction scenes can be described in natural language, but dataset size and richness strongly affect performance. & Connects to VisualSiteDiary and Xiao et al. by starting the captioning-to-reporting line of work. & Captioning gives readable site descriptions, but the approach is mostly static and depends on strong datasets. & Use in Theme 3 and Theme 4 for captioning, data scarcity, and early multimodal direction. \\

Gil and Lee (2024) & PPE monitoring. Core safety paper. & Zero-shot PPE monitoring based on image captioning and visual-language matching. & The method can distinguish PPE-related classes without large task-specific annotation sets. & Extends Nath et al. by reducing the need for bounding-box training data. It also connects to Wang et al. through visual-text safety interpretation. & Zero-shot flexibility is useful, but it may be less stable than supervised detection in complex field images. & Use in Theme 2 as a core example of VLM-adjacent safety monitoring. \\

Chen et al. (2024) & Worker safety query and training. Core safety paper. & Augmented reality, deep learning, and vision-language query system. & The system supports visual safety queries and real-time warnings in an interactive setting. & Connects to Hussain et al. because both move toward assistant-like safety support. & Interaction is a strength, but deployment depends on interface design, real-time performance, and field usability. & Use in Theme 2 to show that safety monitoring is becoming interactive. \\

Wang et al. (2024) & Hazard identification. Core safety paper. & Visual-text semantic similarity between image content and safety rules. & The study links visual site evidence to textual safety requirements instead of stopping at object labels. & Connects Wu et al.'s semantic reasoning with newer VLM-style safety interpretation. & Rule matching adds meaning, but it depends on how well the visual description and rule language align. & Use in Theme 2 to show movement from PPE checks to contextual hazard interpretation. \\

Chan et al. (2025) & Construction safety monitoring. Core VLM paper. & Context-aware VLM agent enriched with domain-specific ontology and prompting. & Domain knowledge improves safety monitoring because generic VLMs lack construction context. & Extends Wu et al. by placing ontology closer to the VLM workflow. Connects to Hussain et al. as an agent-based safety direction. & The approach addresses hallucination and context gaps, but it still depends on prompt quality and domain data. & Use in Theme 1, Theme 2, and Theme 4 for domain adaptation and VLM safety agents. \\

Hussain et al. (2026) & Onsite safety inspection. Core VLM paper. & Vision-language intelligent assistant for safety inspection. & The paper frames safety inspection as an assistant-supported process that combines visual evidence with language-based interpretation. & Connects to Chen et al. and Chan et al. as part of the assistant and agent branch. & Assistant systems are promising, but cross-site reliability and accountability remain open issues. & Use in Theme 2 as a current endpoint of the safety-inspection branch. \\

Bui et al. (2026) & Worker action and emotion understanding. Core benchmark paper. & Evaluation of GPT-4o, Florence 2, and LLaVA-1.5 on construction worker images. & GPT-4o performed best, but all models struggled with close categories and static image limits. & Connects to Sherafat et al. and Jeoung et al. because it shows the need for stronger temporal modeling. & General VLMs can interpret some worker states, but static images are weak for dynamic behavior. & Use in Theme 4 for benchmark evidence and in Theme 3 for the temporal reasoning gap. \\

Jung et al. (2024) & Daily reporting and image retrieval. Core progress paper. & VisualSiteDiary, a detector-free vision-language transformer for photolog captioning. & The model improves automated daily construction reporting and retrieves site images through text descriptions. & Extends Liu et al.'s image captioning work into a more practical reporting pipeline. & Reporting quality improves, but generated text still depends on dataset coverage and evaluation choices. & Use in Theme 3 as a key progress-reporting example. \\

Xiao et al. (2024) & Automated daily reports from video. Core progress paper. & Computer vision combined with ChatGPT-style report generation. & The system links video-based productivity analysis with readable report generation. & Connects to Jung et al. through reporting and to Sherafat et al. through activity tracking. & The method is closer to workflow use, but its evaluation mixes detection accuracy and user judgment. & Use in Theme 3 and Theme 4 for report generation and metric variation. \\

Jeoung et al. (2025) & Equipment task monitoring. Core progress paper. & Zero-shot framework using MLLMs on video frames for spatial, task, and temporal event questions. & The study directly separates spatial, task, and temporal questions for equipment monitoring. & Connects to Bui et al. because both expose the limits of static or short-frame reasoning. It also extends earlier activity-recognition work. & The study is valuable for temporal reasoning, but the field still lacks long-horizon, real-site temporal benchmarks. & Use in Theme 3 as the strongest direct evidence for spatiotemporal task monitoring. \\

Zhang et al. (2025) & Tool and material recognition. Core progress paper. & Training-free few-shot detection using pretrained vision-language methods such as CLIP and DINO-style models. & The study shows that pretrained visual-language models can reduce retraining demands for construction object recognition. & Connects to Liang et al. and Gil and Lee through few-shot and zero-shot recognition. & Few-shot methods reduce data burden, but open-vocabulary recognition still needs reliable prompts and domain terms. & Use in Theme 1 and Theme 4 for open-vocabulary recognition and data-efficiency gaps. \\

\end{longtable}
\end{landscape}

\section{Theme Development and Categorization}

The four themes below organize the matrix into a draft structure for the review. The themes are related, but each one has a different job. Theme 1 explains the technical shift. Theme 2 focuses on safety. Theme 3 focuses on progress and time. Theme 4 names the barriers that cut across the field.

\subsection{Theme 1: From Detection to Semantic Multimodal Reasoning}

\textbf{Theme definition.} Earlier studies often detect visible objects, workers, equipment, PPE, or actions. Newer studies try to interpret meaning. They connect images and videos to captions, questions, rules, prompts, or domain knowledge.

\textbf{Supporting studies.} Zhong et al. (2019), Nath et al. (2020), Wu et al. (2021), Liu et al. (2020), Gil and Lee (2024), Wang et al. (2024), Chan et al. (2025), Hussain et al. (2026), Bui et al. (2026), and Zhang et al. (2025).

\textbf{Synthesis.} The field does not move in one clean step from CV to VLMs. It moves through several layers. Nath et al. show a strong detection baseline for PPE. Wu et al. add semantic reasoning to detection. Liu et al. introduce language through scene captioning. Gil and Lee, Wang et al., Chan et al., Hussain et al., Bui et al., and Zhang et al. show the newer direction: models that connect visual evidence with text, prompts, rules, or domain knowledge.

\textbf{Tension.} Traditional detection can be easier to evaluate. It has clear labels and common metrics. Semantic multimodal reasoning is more useful for real safety and progress decisions, but it is harder to validate. It can be affected by prompt wording, missing domain knowledge, hallucination, and weak construction-specific datasets.

\subsection{Theme 2: Safety Compliance Moves Toward Contextual Hazard Interpretation}

\textbf{Theme definition.} Construction safety monitoring begins with visible compliance checks, such as hardhats and vests. It then expands toward worker behavior, unsafe proximity, safety rules, visual-text matching, ontology-guided reasoning, and assistant-style inspection.

\textbf{Supporting studies.} Nath et al. (2020), Gil and Lee (2024), Chen et al. (2024), Wang et al. (2024), Chan et al. (2025), Hussain et al. (2026), and Bui et al. (2026).

\textbf{Synthesis.} PPE detection is a practical starting point because it is visible and measurable. Nath et al. show the value of real-time detection. Gil and Lee reduce the need for large labeled PPE datasets through zero-shot monitoring. Wang et al. move closer to safety reasoning by matching visual evidence with safety rule language. Chen et al. and Hussain et al. shift the task toward interactive assistance. Chan et al. add construction ontology to support more context-aware decisions.

\textbf{Tension.} Safety compliance is not only a visual question. A worker may be wearing PPE and still be unsafe because of equipment proximity, work sequence, fall risk, or site context. This creates a gap between detection accuracy and practical safety judgment.

\subsection{Theme 3: Progress and Activity Tracking Need Stronger Temporal Reasoning}

\textbf{Theme definition.} Progress monitoring and activity tracking require models to understand changes over time. This includes worker actions, equipment tasks, daily reporting, productivity, and sequence-level site activity.

\textbf{Supporting studies.} Sherafat et al. (2020), Liu et al. (2020), Jung et al. (2024), Xiao et al. (2024), Jeoung et al. (2025), and Bui et al. (2026).

\textbf{Synthesis.} Sherafat et al. show that worker and equipment activity recognition was already a major research area before current VLMs. Liu et al. show that image captioning can describe construction activity scenes. Jung et al. and Xiao et al. make this more practical by connecting images or videos to daily reporting. Jeoung et al. goes further by separating spatial, task, and temporal event questions for equipment monitoring. Bui et al. is important because it shows that static images remain weak for dynamic worker behavior.

\textbf{Tension.} Many papers use images, short clips, or frame-based inputs. Those inputs can support useful descriptions, but they do not fully solve long-horizon reasoning. Construction activity is continuous. The review should therefore separate static visual semantics from true temporal monitoring.

\subsection{Theme 4: Dataset, Benchmark, and Deployment Limits Shape the Field}

\textbf{Theme definition.} The same problems appear across safety and progress studies: limited construction datasets, inconsistent metrics, domain shift, occlusion, site noise, scarce real-site validation, and lack of shared benchmarks.

\textbf{Supporting studies.} Liu et al. (2020), Davila Delgado and Oyedele (2021), Tohidifar et al. (2024), Xin et al. (2024), Bui et al. (2026), Li et al. (2025), Rabbi and Jeelani (2024), Zhong et al. (2019), Chan et al. (2025), and Xiao et al. (2024).

\textbf{Synthesis.} Data is not a side issue. It shapes what the models can do. Liu et al. show that captioning performance depends heavily on dataset size and richness. Bui et al. show that general VLMs struggle with close construction worker categories. Chan et al. identifies the need for construction-specific context and domain knowledge. Xiao et al. shows that evaluation may include both model performance and user judgment. The result is a field with many useful demonstrations, but limited direct comparability.

\textbf{Tension.} Synthetic data, few-shot learning, zero-shot learning, and pretrained VLMs reduce annotation cost. They do not remove the need for reliable field validation. A model that works on curated images may still struggle with occlusion, poor lighting, clutter, camera angle, or site-specific terminology.

\section{Organizational Outline for the Literature Review}

\subsection{I. Introduction}

\begin{enumerate}
\item Open with the practical problem. Construction sites are dynamic work environments where safety and progress information changes quickly.
\item Define the review topic. The review examines VLMs and related multimodal AI systems for safety monitoring and activity tracking on active construction sites.
\item Define spatiotemporal reasoning in plain language. The review uses the term to mean understanding what is happening, where it is happening, and how it changes over time.
\item State the review purpose. The purpose is to map model architectures, data modalities, construction tasks, evaluation practices, deployment barriers, and gaps.
\item State the method. This is a scoping review using PRISMA-ScR guidance because the literature is broad, mixed, and still developing.
\end{enumerate}

\textbf{Guiding questions.}

\begin{itemize}
\item What VLM or multimodal architectures are being used for construction-site safety monitoring and activity tracking?
\item How are visual, textual, temporal, and geometric data combined to support construction safety or progress decisions?
\item What barriers appear when these models are used for real-time or field-based monitoring?
\item How are these systems evaluated, and what metrics or benchmarks are used across studies?
\end{itemize}

\subsection{II. Scoping Review Approach}

This section will explain the search and charting process. It will not read like a full systematic review methods section, but it will keep the review transparent.

\begin{enumerate}
\item Identify the databases and search process, including Web of Science, Scopus, and backward snowballing.
\item Explain inclusion and exclusion criteria. Core studies should focus on active construction sites, safety monitoring, progress tracking, activity recognition, or multimodal construction scene understanding.
\item Explain why baseline papers are included. They show the shift from manual, sensor-based, NLP-only, or CV-only systems toward multimodal reasoning.
\item Describe charting categories. These include model type, data modality, construction task, temporal component, evaluation metric, deployment context, and limitation.
\end{enumerate}

\subsection{III. Theme 1: From Detection to Semantic Multimodal Reasoning}

This body section will explain the technical shift across the literature.

\begin{itemize}
\item Start with baseline CV and detection work, using Zhong et al., Nath et al., and Sherafat et al.
\item Show the transition toward semantic reasoning, using Wu et al. and Liu et al.
\item Then discuss the VLM and multimodal turn, using Gil and Lee, Wang et al., Chan et al., Hussain et al., Bui et al., and Zhang et al.
\item End by stating the main synthesis point. The field is moving from seeing objects to interpreting construction meaning.
\end{itemize}

\subsection{IV. Theme 2: Safety Compliance and Contextual Hazard Interpretation}

This body section will focus on safety monitoring.

\begin{itemize}
\item Begin with PPE and visible compliance detection.
\item Compare zero-shot PPE monitoring with supervised detection.
\item Discuss visual-text safety rule matching and ontology-supported reasoning.
\item Include assistant-style and query-based systems.
\item End with the gap. Practical safety judgment requires context, not only object labels.
\end{itemize}

\subsection{V. Theme 3: Activity Tracking, Progress Monitoring, and Temporal Reasoning}

This body section will focus on progress and time.

\begin{itemize}
\item Begin with activity recognition for workers and equipment.
\item Discuss image captioning and photolog captioning as a bridge into reporting.
\item Compare VisualSiteDiary and video-to-report systems.
\item Use Jeoung et al. to discuss spatial, task, and temporal event questions.
\item End with the gap. Many systems describe frames or short clips, but long-horizon site reasoning remains underdeveloped.
\end{itemize}

\subsection{VI. Theme 4: Datasets, Metrics, and Field Deployment Barriers}

This body section will connect the limits across all themes.

\begin{itemize}
\item Discuss limited construction-specific datasets.
\item Discuss synthetic data and few-shot or zero-shot methods as partial solutions.
\item Compare evaluation metrics across detection, captioning, VQA, report generation, and user studies.
\item Discuss field deployment issues such as occlusion, lighting, clutter, prompt reliability, domain shift, and real-time use.
\item End with the gap. The field needs shared benchmarks and stronger real-site validation.
\end{itemize}

\subsection{VII. Conclusion Placeholder}

The conclusion will summarize the main pattern. Construction AI research is moving from detecting visible site elements toward interpreting safety and progress meaning through language, context, and multimodal reasoning. The strongest future directions are construction-specific VLM benchmarks, better temporal datasets, transparent field validation, and systems that connect visual observations to safety rules and work progress in ways practitioners can trust.

\section{Submission Note}

This document contains all three required components in one formatted file: the synthesis matrix, the theme development section, and the organizational outline. If the course platform requires separate uploads, the matrix and theme sections can be exported as the working document, and the outline section can be submitted as the separate outline document.

\end{document}
